% Partie : counterfactuals
% Chapitre : minimal-change-semantics
% Section : sphere-models

\documentclass[../../../include/open-logic-section]{subfiles}

\begin{document}

\olfileid[fr]{cnt}{min}{sph}

\olsection{Les modèles à sphères}

Pour donner une sémantique formelle des conditionnels contrefactuels,
on peut transformer l'analyse informelle de Lewis en structure
mathématique. Les sphères autour d'un monde~$w$ sont alors des
ensembles de mondes. Puisqu'elles sont emboîtées, ces ensembles
autour de~$w$ doivent être totalement ordonnés par l'inclusion.

\begin{defn}
  Un \emph{modèle à sphères} est un triplet~$\mModel{M} = \tuple{W, O, V}$
  où $W$ est un ensemble non vide de mondes, $V\colon \PVar \to \Pow{W}$
  est une valuation, et $O\colon W \to \Pow{\Pow{W}}$ associe à chaque
  monde~$w$ un \emph{système de sphères}~$O_w$. Pour chaque~$w$,
  $O_w$ est un ensemble d'ensembles de mondes qui doit satisfaire les
  conditions suivantes :
  \begin{enumerate}
  \item $O_w$ est \emph{centré} sur~$w$ : $\{w\} \in O_w$ ; aucune
    sphère du système n'est vide.
  \item $O_w$ est \emph{emboîté} : si $S_1$, $S_2 \in O_w$, alors $S_1
    \subseteq S_2$ ou $S_2 \subseteq S_1$ ; autrement dit, $O_w$ est
    totalement ordonné par~$\subseteq$.
  \item $O_w$ est stable par réunion de toute sous-famille non vide.
  \item $O_w$ est stable par intersection de toute sous-famille non vide.
  \end{enumerate}
\end{defn}

L'idée qui sous-tend $O_w$ est que les mondes « autour » de~$w$ sont
répartis selon leur distance à~$w$. La sphère la plus intérieure ne
contient que~$w$, c'est-à-dire l'ensemble~$\{w\}$ : $w$ est plus proche
de~$w$ que n'importe quel autre monde figurant dans les sphères. Si
$S \subsetneq S'$,
les mondes de $S' \setminus S$ sont plus éloignés de~$w$ que ceux de~$S$ :
$S' \setminus S$ forme la « couche » comprise entre~$S$ et les mondes
extérieurs à~$S'$. Il faut donc concevoir les sphères comme contenant
tous les mondes situés à l'intérieur de leur frontière, et non
simplement les mondes d'une couche donnée.

\begin{figure}
\begin{center}
\begin{tikzpicture}[layerwidth=1.5,scale=.6]\tiny
  \spheresystem{3}
  \propositionintersect{0}{3}{60}{6}
  \path (0,0) node[world] {$w$};
  \spherepos{-30}{2}{node[world] {$w_2$}}
  \spherepos{220}{2}{node[world] {$w_3$}}
  \spherepos{90}{2}{node[world] {$w_1$}}
  \spherepos[shift={(.1,0)}]{10}{3}{node[world] {$w_5$}}
  \spherepos[shift={(.1,0)}]{-10}{3}{node[world] {$w_6$}}
  \spherepos{225}{3}{node[world] {$w_4$}}
  \spherepos{20}{4}{node[world] {$w_7$}}
  \path (5,0) node[left] {$p$};
\end{tikzpicture}
\caption{Schéma d'un modèle à sphères}
\ollabel{fig:sphere-model}
\end{center}
\end{figure}

Le schéma de la \olref{fig:sphere-model} représente le modèle à
sphères où $W = \{w, w_1, \dots, w_7\}$ et $V(p) = \{w_5, w_6,
w_7\}$. La sphère la plus intérieure est $S_1 = \{w\}$. Les mondes
les plus proches de~$w$, hormis le monde central lui-même, sont
$w_1, w_2, w_3$ ; la sphère suivante est donc
$S_2 = \{w, w_1, w_2, w_3\}$. Plus loin se
trouvent $w_4$, $w_5$ et~$w_6$ ; la sphère extérieure est donc
$S_3 = \{w, w_1, \dots, w_6\}$. Le système de sphères autour de~$w$
est $O_w = \{S_1, S_2, S_3\}$. Le monde~$w_7$ n'appartient à aucune
sphère autour de~$w$. Les mondes les plus proches où~$p$ est vraie
sont $w_5$ et~$w_6$ ; la plus petite sphère contenant des mondes
où~$p$ est vraie est donc~$S_3$.

Pour définir la satisfaction d'une formule~$!A$ au monde~$w$ dans un
modèle à sphères~$\mModel M$, soit $\mSat{M}{!A}[w]$, nous ajoutons à
la définition donnée pour les !!{formula}s modales une clause pour
$!B \cif !C$ :

\begin{defn}
  $\mSat{M}{!B \cif !C}[w]$ si et seulement si l'une des deux
  conditions suivantes est remplie :
  \begin{enumerate}
  \item\ollabel{sphere-vac} Pour tout $u \in \bigcup O_w$,
    $\mSat/{M}{!B}[u]$, ou
  \item\ollabel{sphere-nonvac} pour une certaine sphère $S \in O_w$,
    \begin{enumerate}
      \item $\mSat{M}{!B}[u]$ pour au moins un $u \in S$, et
      \item pour tout $v \in S$, soit $\mSat/{M}{!B}[v]$, soit
        $\mSat{M}{!C}[v]$.
    \end{enumerate}
  \end{enumerate}
\end{defn}

Selon cette définition, $\mSat{M}{!B \cif !C}[w]$ si et seulement si
l'antécédent~$!B$ est faux partout dans les sphères autour de~$w$,
ou s'il existe une sphère~$S$ dans laquelle $!B$ est vraie et où le
conditionnel matériel $!B \lif !C$ est vrai dans tous les mondes de
cette sphère « admettant~$!B$ ». La définition n'exige pas, contrairement
à ce que pourrait suggérer l'explication intuitive, que~$S$ soit la
plus petite sphère \emph{admettant~$!B$}. Toutefois, si la condition
de l'\olref{sphere-nonvac} est remplie pour une sphère~$S$, elle
l'est aussi pour toute sphère contenue dans~$S$ \emph{qui admet
  aussi l'antécédent} et donc, si elle existe, pour la plus petite sphère de
ce type.

Une plus petite sphère admettant~$!B$ peut ne pas exister.
Prenons $W=\{w\}\cup\{u_1,u_2,\ldots\}$ et les sphères
$S_i=\{w\}\cup\{u_j:j\ge i\}$, pour $i\ge 1$. Au monde~$w$,
le système est constitué des~$S_i$ et de~$\{w\}$. Au monde~$u_i$,
prenons le système réduit à la sphère~$\{u_i\}$. Au centre~$w$,
la réunion d'une famille contenant des~$S_i$ est le~$S_i$ de plus
petit indice dans la famille. Si la famille ne contient que des
$S_i$, son intersection est le~$S_i$ de plus grand indice
lorsqu'il existe, et $\{w\}$ sinon. La présence de~$\{w\}$ rend
l'intersection égale à~$\{w\}$. Le système satisfait donc les
conditions de fermeture.

Supposons que $!B$ soit vrai à tous les~$u_i$ et faux en~$w$.
Toutes les sphères $S_i$ admettent~$!B$, et chacune contient
une sphère plus petite qui l'admet encore. Leur intersection
$\{w\}$ n'admet pas~$!B$. Dans cet exemple, le contrefactuel
est donc vrai en~$w$ si et seulement si $!C$ est vrai en tous
les~$u_j$ à partir d'un certain indice. Plus généralement, on
peut tester une suite décroissante de sphères admettant~$!B$
lorsque toute sphère admettant~$!B$ contient au moins un terme
de la suite. Tout témoin donné par la définition contient alors
un terme de la suite qui est aussi un témoin.

\end{document}
